%%%%%%%%%%%%%%%%%%%%%%%%%%%%%%%%%%%%%%%%%%%%%%%%%%%%%%%%%%%%%%%
%
% Welcome to Overleaf --- just edit your LaTeX on the left,
% and we'll compile it for you on the right. If you open the
% 'Share' menu, you can invite other users to edit at the same
% time. See www.overleaf.com/learn for more info. Enjoy!
%
%%%%%%%%%%%%%%%%%%%%%%%%%%%%%%%%%%%%%%%%%%%%%%%%%%%%%%%%%%%%%%%
\documentclass{beamer}
\usepackage{tikz}
\usetheme{Copenhagen}
\usepackage{booktabs}
\usecolortheme{default}
\setbeamertemplate{headline}{}
\title[Online Mirror Descent, Regret Bounds and Applications] %optional
{Online Convex Optimisation - Part II}
\subtitle{Online Mirror Descent, Regret Bounds and Applications}
\author[CS6216 Group 8a] % (optional)
{Teoh Tze Tzun\inst{1} \and Tan Kai Min, Russell\inst{1}}

\institute[VFU] % (optional)
{
  \inst{1}%
  School of Computing\\
  National University of Singapore
}

\date[VLC 2021] % (optional)
{CS6216 Group 8a}


%End of title page configuration block
%------------------------------------------------------------
%The next block of commands puts the table of contents at the 
%beginning of each section and highlights the current section:

\AtBeginSection[]
{
  \begin{frame}
    \frametitle{Table of Contents}
    \tableofcontents[currentsection]
  \end{frame}
}
%------------------------------------------------------------
\begin{document}
\frame{\titlepage}
%---------------------------------------------------------
\section{Online Convex Optimisation}
\begin{frame}{What is Online Convex Optimisation (OCO)}
\pause
\begin{columns}[t] % Aligns content at the top
        \begin{column}{0.48\textwidth}
            \textbf{Convex Optimisation (CO)}
            \begin{align*}\\
                \text{minimise } &f(x)\\
                \text{subject to }&g^{(i)}(x) \geq 0\\
                &h_j(x) =0
            \end{align*}
        where $f$ is convex and the feasible set due to the constraints are also convex
        \end{column}\pause       
        \begin{column}{0.48\textwidth}
            \textbf{Online Convex Optimisation (OCO)}
            \begin{align*}
                \text{Input }&: \text{Convex Set }S\\
                \text{For }t&=1\ldots T:\\
                &\text{predict } x^{(t)} \in S\\
                &\text{receive loss function }f^{(t)}: S \to \mathbb R\\
                &\text{learn loss }f^{(t)}(x^{(t)})
            \end{align*}
            where $f^{(t)}$ is a convex function.
        \end{column}
\end{columns}
\vspace{0.5cm}
\pause
Their objectives are also different. It is clear that in CO that we want to minimise our objective function, but in OCO we minimise something called the \textbf{regret}.
\end{frame}

\begin{frame}{Regret}
\begin{block}{Definition of Regret}
Regret with respect to a fixed choice $x$ is defined as
    \begin{equation}
    \text{Regret}^{(t)}(x) = \underbrace{\sum_{t=1}^T f^{(t)}(w^{(t)})}_{\text{Algorithm's loss}} - \underbrace{\sum_{t=1}^Tf^{(t)}(x)}_{\text{Fixed choice total loss}}
\end{equation}
\end{block}
\pause
Intuitively, regret is the extra total loss an algorithm suffers over time compared to the \textit{single best fixed choice it could have made if it had known all the data from the beginning}.
\end{frame}

\section{Recap from Part I}
\begin{frame}{Follow the Regularised Leader (FoReL)}
In Week 7, the group gave the following frameworks for OCO:\pause
\begin{itemize}
    \item Follow-The-Leader (FTL): 
    select the next iterate my minimising the loss on all past rounds.
    \begin{equation}
         x^{(t)} = \arg \min_{x \in S} \left\{\sum_{i = 1}^{t-1} f^{(i)}(x)\right\}\qquad\forall\,t
    \end{equation}\pause
    \item Follow-The-Regularised-Leader (FoReL):
    extends FTL by adding a regularisation term $h: S \to \mathbb R$ for stability
    \begin{equation}
        x^{(t)} = \arg \min_{x \in S} \left\{\sum_{i = 1}^{t-1} f^{(i)}(x) + h(x)\right\}\qquad\forall\,t
    \end{equation}
\end{itemize}
\end{frame}

\begin{frame}{Regret Bound Results for FoReL}
Suppose we assume the following:
\begin{itemize}
    \item Regulariser $h$ is $\boldsymbol{\sigma}$\textbf{-strongly convex} w.r.t $||\cdot||$ norm.
    i.e for any $x' \in S$, we have that $\forall z \in \partial h(x'),\;\forall x \in S,$\begin{equation*}
         h(x) \geq h(x') + \langle z, x - x' \rangle + \frac{\sigma}{2} || x - x' ||^2
    \end{equation*}
\pause
    \item $f^{(t)}$ is $\mathbf{L_t}$\textbf{-Lipschitz} i.e for all $x, x' \in S$ we have $$|f^{(t)}(x) - f^{(t)}(x')| \leq L_t||x-x'||$$
\end{itemize}
We have the following bounds:
        \begin{equation}\label{eqn:forel_regre_with_min}
        \begin{aligned}
    \text{Regret}_T(x) &\leq h(x) - h(x^{(1)}) + \frac{TL^2}{\sigma} \leq h(x) - \min_{v \in S} h(v) + \frac{TL^2}{\sigma}
    \end{aligned}
\end{equation} 
where $L$ is such that $\frac{1}{T}\sum_{t=1}^T L_t^2 \leq L^2$.
\end{frame}

\begin{frame}{Proof Sketch of Regret Bound}
We first have a lemma from Week 7: \begin{lemma}\label{lemma:iterates_forel}
    Let $x^{(1)}, x^{(2)}, \ldots$ be the iterates due to FoReL using the $\sigma$-strongly convex regulariser $h$. Then for all $x \in S$, we have that \begin{equation}\label{eqn:lemma_forel}
        \sum_{t=1}^T (f^{(t)}(x^{(t)}) - f^{(t)}(x)) \leq h(x) - h(x_1) + \sum_{t=1}^T (f^{(t)}(x^{(t)}) - f^{(t)}(x^{(t+1)}))
    \end{equation}
\end{lemma}\pause
Furthermore, it is also shown that if for all $t$, if $f^{(t)}$ is $L^{(t)}$-Lipschitz with respect to $||\cdot||$, then we further have that 
\begin{equation}\label{eqn:lemma_forel_lipschitz}
    f^{(t)}(x^{(t)}) - f^{(t)}(x^{(t+1)}) \leq L_t ||x^{(t)} - x^{(t+1)}|| \leq \frac{L_t^2}{\sigma}
\end{equation} \pause
    Combining these 2 results, we obtain the expression of the regret bound.
\end{frame}

\begin{frame}{The Problem with FoReL}
    \begin{itemize}
        \item Update/Learning Rule: $$x^{(t)} = \arg \min_{x \in S} \left\{\sum_{i = 1}^{t-1} f^{(i)}(x) + h(x)\right\}\qquad\forall\,t$$ requires solving an optimisation problem at every iteration.
        \item Can be computationally expensive and complex
    \end{itemize}\pause
    \begin{block}{Motivating Question}
    Can we achieve a similar regret bound with a \textit{simpler} update rule?
    \end{block}
    \pause
    Online Mirror Descent gives a much simpler update rule with similar regret bounds!
\end{frame}
\section{Online Mirror Descent}
\begin{frame}{What is Classical Mirror Descent}
\begin{itemize}
    \item Consider the standard gradient descent update rule: $$x^{(t+1)} = x^{(t)} - \eta \nabla f(x^{(t)})$$\pause
    \item Underlying assumption: The iterates ($x^{(1)}, \ldots, x^{(T)}$) are assumed to be in the same space as the gradients ($\nabla f(x^{(1)}), \ldots, \nabla f(x^{(T-1)})$) - which is not true in general! \pause
    \item E.g what happens if your search space is the $d$-dimension probability simplex $\Delta_d$?
\end{itemize}
\begin{block}{Main idea of Mirror Descent}
    Assume the iterate-space (primal space) and the gradient-space (dual space) are spaces of 2 distinct geometries, $\mathcal{X}, \theta$ respectively
\end{block}
\end{frame}

\begin{frame}{The Primal-Dual Space Map}
We therefore need a way to ``map" from the primal space $\mathcal{X}$ to the dual space $\theta$. This map also has to be \textbf{invertible}, as we want to obtain our next iterate in the primal space too. To define this map:
\begin{itemize}
    \item Fix a norm $||\cdot||$
    \item Choose some differentiable, convex function $h: \mathbb{R}^n \to \mathbb{R}$ that is $\sigma$-strongly convex.\pause
    \item Use $h$ to define a regulariser called the \textit{Bregman Divergence}
\end{itemize}
\end{frame}

\begin{frame}{Bregman Divergence, $D_h$}
    \begin{block}{Bregman Divergence}
        The Bregman divergence with respect to a strictly convex function $h$ (denoted by $D_h$), of two points, $x$ and $x'$ is defined as follows:
        \begin{equation}
            D_h(x || x') := h(x) - h(x') - \langle\nabla h(x'),  x - x' \rangle
        \end{equation}
    \end{block}
    \pause
    \begin{columns}[c]
        \begin{column}{0.5\textwidth}
            \begin{figure}
                \centering
                \includegraphics[width=\linewidth]{slides//figures/bregmandiv.png}
                \caption{Enter Caption}
                \label{fig:placeholder}
            \end{figure}
        \end{column}
        \pause       
        \begin{column}{0.48\textwidth}
            \centering 
            Intuitively, it is the \textit{distance} from the function $h$ to its linear approximation.
        \end{column}
    \end{columns}
\end{frame}

\begin{frame}{The Mirror Map}
    Turns out $\nabla h$ is the map that we want. To show that $\nabla h$ is the map:
    \begin{itemize}
        \item Start off with a single proximal gradient step
        \begin{equation}\label{eq:proximal_gradient_update}
        x^{(t+1)} = \arg \min_{x} \left\{\eta \langle \nabla f(x^{(t)}), x\rangle + R(x)\right\} 
        \end{equation} where $R$ is some regularisation function. 
    \end{itemize}

    
\end{frame}

\begin{frame}{The Mirror Map}
    Turns out $\nabla h$ is the map that we want. To show that $\nabla h$ is the map:
    \begin{itemize}
        \item Start off with a single proximal gradient step
        \begin{equation}\label{eq:proximal_gradient_update}
        x^{(t+1)} = \arg \min_{x} \left\{\eta \langle \nabla f(x^{(t)}), x\rangle + R(x)\right\} 
        \end{equation} where $R$ is some regularisation function. 
    \pause
    \item Set $R(x) = D_h(x||x^{(t)})$ based on the $h$ we have chosen
    \pause
    \item Solve the minimisation problem of Equation~\ref{eq:proximal_gradient_update} to get the following update:
    \begin{equation}
        \nabla h(x^{(t+1)}) \gets \nabla h(x^{(t)}) - \eta \nabla f(x^{(t)})
    \end{equation}
    \end{itemize}
    \pause
    Notice how instead of $x^{(t)}$, we have $\nabla h(x^{(t)})$. We denote this ``mapped" variable as $\theta^{(t)} = \nabla h(x^{(t)})$
\end{frame}


\begin{frame}{Mirror Descent Algorithm}
    \begin{columns}[t] % Aligns content at the top
        \begin{column}{0.48\textwidth}
            \textbf{Gradient Descent}
            \begin{align*}
                x^{(t+1)} &= x^{(t)} - \eta \nabla f(x^{(t)})
            \end{align*}
        \end{column}\pause       
        \begin{column}{0.48\textwidth}
            \textbf{Mirror Descent}
            \begin{align*}
                \theta^{(t)}&= \nabla h(x^{(t)})\\
                \theta^{(t+1)} &\gets \theta^{(t)} - \eta \nabla f(x^{(t)})\\
                x^{(t+1)}&=\Pi((\nabla h)^{-1}(\theta^{(t+1)}))
            \end{align*}
        \end{column}
\end{columns}
\vspace{0.5cm}\pause
Notice for Mirror Descent, we do the gradient step in the dual space instead of the primal space. Also note the following:
\begin{enumerate}
    \item $h$ is strongly-convex $\implies \nabla h$ is invertible
    \item The inverse map might not give a feasible $x$. A projection $\Pi$ might be needed to ensure the next iterate is feasible
\end{enumerate}
\end{frame}

\begin{frame}{Mirror Descent Algorithm}
    \begin{figure}
        \centering
        \includegraphics[width=\linewidth]{slides/figures/mirror_descent.png}
        \caption{Mirror Descent}
        \label{fig:mirror_descent}
    \end{figure}
\end{frame}

\begin{frame}{Online Mirror Descent (OMD)}
    Mirror Descent but with the following differences:
    \begin{enumerate}
    \item \textbf{Objective Function} The objective function for online mirror descent changes at every step. At step $t$, the algorithm must commit to a prediction $x^{(t)}$ before the ``environment" reveals the specific loss function $f^{(t)}$ for that iteration.
    \item \textbf{Gradient} In online mirror descent, we use the \textit{sub-gradient} of the newly revealed function $f^{(t)}$ for that iteration. We denote the sub-gradient as $z^{(t)} \in \partial f^{(t)}(x^{(t)})$. This implies that the update rule is now \begin{equation}
        \theta^{(t+1)}=  \theta^{(t)} - \eta \cdot z^{(t)}\qquad z^{(t)} \in \partial f^{(t)}(x^{(t)})
    \end{equation}
\end{enumerate}
\end{frame}

\begin{frame}{A note on subgradients}
    A sub-gradient of a function is a more generalised form of a gradient. More specifically, we define sub-gradients as follows:
    \begin{block}{Sub-gradients}\label{defn:subgradient}
    A value $z$ is a sub-gradient of $f$ at $x$ if for all $x' \in S$\begin{equation}
        f(x') \geq f(x) + \langle x' - x, z\rangle
    \end{equation}
    The set of sub-gradients of $f$ at $x'$ is denoted by $\partial f(x')$. 
\end{block}\pause
Furthermore, if $f$ is differentiable at $x'$, then $\partial f(x') = \{\nabla f(x')\}$ (i.e $\partial f(x^*)$ is a singleton and its only element is the gradient of $f$ at $x'$).
\end{frame}

\begin{frame}{A Generic Regret Bound}
\begin{block}{Generic Regret Bound for OMD Algorithms}
Suppose we select a $\frac{1}{\eta}$-strongly convex regulariser $h$ and let $f_t$ be linear functions that are $L_t$-Lipschitz. Then we have the following generic regret bound
    \begin{equation*}
        \text{Regret}_T(x) \leq h(x) - \min_{v\in S}h(v) + \eta \sum_{t=1}^T ||z^{(t)}||_*^2
    \end{equation*}
\end{block}\pause
The old framework FoReL can help us achieve that bound!
\end{frame}

\begin{frame}{Linearisation Trick - FoReL and OMD}
    Recall the mirror descent step:
    \begin{align}
        x^{(t+1)} \gets (\nabla h)^{-1}(\theta^{(t+1)}) := \mathbf g(\theta^{(t+1)})
    \end{align}
    We can re-write $g(\theta^{(t+1)}) = \arg\max_{x \in S} \{\langle \theta^{(t+1)}, x\rangle - h(x)\}$ 
    
    To show this:
    \begin{itemize}
        \item Start with Proximal Gradient Step: $x^{(t+1)} = \arg \min_{x} \left\{\eta \langle \nabla f^{(t)}(x^{(t)}), x\rangle + D_h(x || x^{(t)})\right\}$
        \item Expand expression of Bregman divergence and remove constants to get $\arg\min_x \{ \eta \langle \nabla f^{(t)}(x^{(t)}), x \rangle + h(x) - \langle \nabla h(x^{(t)}), x \rangle \}$
        \item Factor out $x$ term and substitute the update rule $\theta^{(t+1)} =\nabla h(x^{(t)}) - \eta \nabla f^{(t)}(x^{(t)})$
    \end{itemize}
\end{frame}

\begin{frame}{Linearisation Trick - FoReL and OMD}
    Suppose we apply FoReL on a sequence of linear functions $f^{(t)}(x) = \langle x, z^{(t)} \rangle$ with some regulariser $h(x)$. \pause
    \begin{align*}
    x^{(t+1)} &= \arg\min_{x \in S} \left\{\sum_{i=1}^t f^{(i)}(x) + h(x) \right\} \\
    &= \arg\min_{x \in S} \left\{\sum_{i=1}^t \langle x, z^{(t)} \rangle  + h(x)\right\} \\
    &= \arg\max_{x \in S} \left\{\left\langle x, \underbrace{-\sum_{i=1}^t z^{(i)}}_{\theta^{(t+1)}} \right\rangle - h(x)\right\}
\end{align*}\pause
We get $g(\theta^{(t+1)})$ from OMD alongside $\theta^{(t+1)} \gets\theta^{(t)} - z^{(t)}.$
\end{frame}

\begin{frame}{Linearisation Trick - FoReL and OMD}
    \begin{block}{}
        OMD $\iff$ FoReL on sequence of linear functions
    \end{block}
    With this relationship, we can borrow regret bounds results from FoReL and re-use them in OMD.\pause
    \begin{itemize}
        \item If $f^{(t)}$ is non-linear but still convex, pick $z^{(t)} \in \partial f(x^{(t)})$\pause
    \end{itemize}
    \begin{equation}\label{eqn:linearisation_trick}
        \underbrace{f^{(t)}(x^{(t)}) - f^{(t)}(x)}_{\text{Regret for non-linear $f^{(t)}$}} \leq \langle x^{(t)} - x, z^{(t)} \rangle = \underbrace{\langle x^{(t)}, z^{(t)}\rangle - \langle x, z^{(t)} \rangle}_{\text{Regret for linear functions}}
    \end{equation}
    The inequality is due to the definition of sub-gradient!
    \begin{block}{Main Takeaway of Linearisation Trick}
        Regret with respect to non-linear functions can still be upper bounded by that of linear functions.
    \end{block}
\end{frame}

\begin{frame}{Proof of Generic Regret Bound}
    \begin{itemize}
        \item Recall OMD $\iff$ FoReL on sequence of linear functions, we use Equation~\ref{eqn:lemma_forel}
    \end{itemize}
    $\text{Regret}_T(x) \leq h(x) - h(x^{(1)}) + \sum_{t=1}^T (f^{(t)}(x^{(t)}) - f^{(t)}(x^{(t+1)}))$
    \begin{alertblock}{Very Important Note on $f^{(t)}$}
        $f^{(t)}$ is w.r.t to FoReL, not the loss function of OMD!
    \end{alertblock}
    
    \begin{itemize}
        \item By Equation~\ref{eqn:lemma_forel_lipschitz}, $\sum_{t=1}^T (f^{(t)}(x^{(t)}) - f^{(t)}(x^{(t+1)})) \leq \sum_{t=1}^T \frac{L_t^2}{\sigma}$
        \item Sub $\sigma= \frac{1}{\eta}$ and substitute $f^{(t)}$ as linear functions instead. Since $f^{(t)}(x) = \langle \theta^{(t)}, x\rangle \implies L_t = ||z_t||_*$
        \item Last step: $\min_{v \in S} h(v) \leq h(x^{(1)})$ yields the final regret bound.
    \end{itemize}
\end{frame}

\begin{frame}{Variants of OMD}
    \begin{itemize}
        \item Normalised Exponentiated Gradient (EG) - $h(x) = \frac{1}{\eta}\sum_{i=1}^d x_i \log x_i$ and norm: $||\cdot||_1$, $S = \Delta_d$
        \item Online Gradient Descent With Lazy Projections - $h(x) = \frac{1}{2\eta}||x||_2^2$ and norm: $||\cdot||_2$
        \item $p$-norm Algorithm ($p\geq 2$) - \newline$h(x) = \frac{1}{2\eta(q-1)}||x||_q^2$ and norm: $\ell_q$ where $\frac{1}{p} + \frac{1}{q} = 1$
        \item Many more, but we will focus on the above 3.
    \end{itemize}
    Note: All the $h$ above are $\frac{1}{\eta}$-strongly convex functions
\end{frame}

\begin{frame}{Normalised Exponentiated Gradient (EG)}
\begin{itemize}
    \item $h(x) = \frac{1}{\eta}\sum_{i=1}^d x_i \log x_i$ (Entropic Regulariser)
    \item Norm: $||\cdot||_1$
    \item $S = \Delta_d$
\end{itemize}\pause

Recall that we derived $g$ relating $\theta^{(t+1)}$ with $x^{(t+1)}$ \pause
\begin{equation*}
    x^{(t+1)} = g(\theta^{(t+1)}) = \arg\max_{x \in \Delta_d} \left\{\sum_{i=1}^d \theta_i^{(t+1)}x_i - \frac{1}{\eta}\sum_{i=1}^d x_i \log x_i\right\}
\end{equation*} which is solvable by Lagrange multipliers (see notes) and it yields the following:
\begin{equation}
    x_i^{(t+1)} = \frac{\exp(\eta \theta_i^{(t+1)})}{\sum_{j=1}^d \exp(\eta \theta_j^{(t+1)})} \hspace{3cm}\text{(Softmax)}
\end{equation}    
\end{frame}

\begin{frame}{Normalised EG Algorithm}
    \begin{enumerate}
    \item Set $w_1 = (1/d, 1/d, \ldots, 1/d)$
    \item For all $i$: $x_i^{(t+1)} =\frac{x_i^{(t)} \exp(-\eta z_i^{(t)})}{\sum_{j=1}^d x_j^{(t)} \exp(-\eta z_j^{(t)})} $ where $z^{(t)} \in \partial f^{(t)}(x^{(t)})$
\end{enumerate}
\pause
\begin{block}{Normalised EG Regret Bound}
    Suppose that $f^{(1)}, \ldots, f^{(T)}$ are a sequence of convex functions in which $f^{(t)}$ is $L_t$-Lipschitz with respect to $||\cdot||_1$. Then
    \begin{equation*}
        \text{Regret}_T(\Delta_d) \leq \frac{\log d}{\eta} + \eta TL^2
    \end{equation*}
    where $L$ is such that $\frac{1}{T}\sum_{t=1}^TL_t^2 \leq L^2$
\end{block}
\end{frame}

\begin{frame}{Regret Bound of Normalised EG}
    We can re-use results from FoReL again!
    \begin{align*}
        \text{Regret}_T(x) &\leq h(x) - \min_{v \in \Delta_d} h(v) + \sum_{t=1}^T \frac{L_t^2}{\sigma}\\
        &\leq h(x) - \min_{v \in \Delta_d} h(v) + \eta TL^2 \qquad \left(\sigma = \frac{1}{\eta}\wedge \frac{1}{T}\sum_{t=1}^TL_t^2 \leq L^2\right)
    \end{align*}\pause
    Observe $h(x)$ is actually negative entropy, which is minimised when $x_i = 1/d$ for all $i$ (Uniform distribution). Hence
    \begin{align*}
        h(x) - \underbrace{\frac{1}{\eta}\sum_{i=1}^d \frac{1}{d} \log \frac{1}{d}}_{\min_{v \in \Delta_d} h(v)} + \eta TL^2= h(x) + \frac{\log d}{\eta}  + \eta TL^2
    \end{align*}
\end{frame}

\begin{frame}{Regret Bound of Normalised EG}
    Final step is that $h(x)$ is maximised at one of the vertices of the simplex (i.e standard basis vector of simplex), hence $h(x) \leq \frac{1}{\eta} \log 1 = 0$ which yields an expression independent of $x$!\pause

    \begin{equation}\label{eqn:neg_regret}
    \boxed{\text{Regret}_T(\Delta_d) \leq \frac{\log d}{\eta} + \eta TL^2}
\end{equation}
In particular, if $\eta = \frac{\sqrt{\log d}}{L \sqrt{2T}}$, then we have $$\text{Regret}_T(\Delta_d) \leq L\sqrt{2T \log d}$$.
\end{frame}

\begin{frame}{Online Gradient Descent with Lazy Projections}
    \begin{itemize}
    \item $h(x) = \frac{1}{2\eta}||x||_2^2$ (Euclidean Regulariser)
    \item Norm: $||\cdot||_2$
    \item $S = \Delta_d$
\end{itemize}\pause
Again, using the definition of $g$, we derive some update rule:\pause
\begin{align*}
    g(\theta) &= \arg\max_{x \in S} \left\{\langle x, \theta\rangle - \frac{1}{2\eta}||x||_2^2\right\} = \arg\max_{x \in S} \left\{\langle x, \eta\theta\rangle - \frac{1}{2}||x||_2^2\right\} \\
    &= \arg\min_{x \in S} \left\{\frac{1}{2}||x||_2^2 - \langle x, \eta\theta\rangle\right\} \\
    &= \arg\min_{x \in S} \left\{\frac{1}{2}||x||_2^2 - \langle x, \eta\theta\rangle + \frac{1}{2}||\eta\theta||_2^2\right\} \hspace{0.3cm}(\because \text{complete the square)}\\
    &= \arg\min_{x \in S} \frac{1}{2}||x - \eta\theta||_2^2 =\arg\min_{x \in S} ||x - \eta\theta||_2 
\end{align*}
\end{frame}

\begin{frame}{OGD with Lazy Projections Algorithm}
    The update rule derived in the previous slide yields the following algorithm:\pause
    \begin{enumerate}
    \item Set $\theta^{(1)} = 0$ (zero or the zero-vector)
    \item for $t = 1 \ldots T$, \begin{enumerate}
        \item $x^{(t)} = \arg\min_{x \in S} ||x - \eta \theta^{(t)}||_2$
        \item $\theta^{(t+1)} \gets \theta^{(t)} - z^{(t)}$ where $z^{(t)} \in \partial f^{(t)}(x^{(t)})$
    \end{enumerate}
\end{enumerate}
Why Lazy?\begin{itemize}
    \item Usual Steps: Gradient Step $\to$ Project to $S$ $\to$ Use projected point as next iterate
    \item This Algorithm: Gradient Step on un-projected point $\to$ Projects it to next iterate
\end{itemize}
It is lazy in the sense that it accumulates gradient updates in an unconstrained space and delays the projection step until a feasible output is strictly required.
\end{frame}

\begin{frame}{Regret Bound for OGD with Lazy Projection}
    \begin{block}{OGD with Lazy Projection Regret Bound}
    Suppose that $f^{(1)}, \ldots, f^{(T)}$ are a sequence of convex functions in which $f^{(t)}$ is $L_t$-Lipschitz with respect to $||\cdot||_2$. Then
    \begin{equation*}
        \text{Regret}_T(x) \leq \frac{1}{2\eta}||x||_2^2 + \eta TL^2
    \end{equation*}
    where $L$ is such that $\frac{1}{T}\sum_{t=1}^TL_t^2 \leq L^2$
\end{block}
\pause
Derivation is simple as it follows the same initial steps as above except the following:
\begin{itemize}
    \item $\min_{v \in S} h(v) = \min_{v \in S}\frac{1}{2\eta}||v||_2^2 \geq 0$ now
    \item We cannot simply upper bound $h(x)$ with some expression independent of $x$.
\end{itemize}
\end{frame}

\begin{frame}{Regret Bound for OGD with Lazy Projection}
    \begin{block}{OGD with Lazy Projection Regret Bound}
    Suppose that $f^{(1)}, \ldots, f^{(T)}$ are a sequence of convex functions in which $f^{(t)}$ is $L_t$-Lipschitz with respect to $||\cdot||_2$. Then
    \begin{equation*}
        \text{Regret}_T(x) \leq \frac{1}{2\eta}||x||_2^2 + \eta TL^2
    \end{equation*}
    where $L$ is such that $\frac{1}{T}\sum_{t=1}^TL_t^2 \leq L^2$
\end{block}
\pause
In particular, if there exists $B \geq \max_{x \in S}||x||_2$, and $\eta = \frac{B}{L\sqrt{2T}}$, we have a new regret bound for the Online Gradient Descent with Lazy Projection algorithm with 
$$\text{Regret}_T(S) \leq BL \sqrt{2T}$$.
\end{frame}

\begin{frame}{$p$-norm Algorithm}
The $p$-norm algorithm ($p \geq 2$ is the parameter) extends the OMD framework by substituting the Euclidean or entropic regulariser with a generalised $q$-norm regulariser.
\begin{itemize}
    \item $h(x) = \frac{1}{2\eta (q - 1)}||x||_q^2$
    \item Norm: $||\cdot||_q$ where $\frac{1}{p} + \frac{1}{q} = 1$. So $p \geq 2 \implies q \in (1, 2]$ 
    \item $S = \mathbb R^d$
\end{itemize}\pause
By using $h(x)$, we arrive at a rather complicated expression of $g(\theta)$ which we will not derive here.
\begin{equation}\label{eqn:p_norm_g}
    g_i(\theta) = \eta \frac{\text{sgn}(\theta_i)  |\theta_i|^{p-1}}{||\theta||_p^{p-2}}
\end{equation} where $p \geq 2$ and $||\theta||_p := \left(\sum_{i=1}^d |\theta_i|^p\right)^{1/p}$. 
\end{frame}

\begin{frame}{Regret Bound for $p$-norm algorithm}
    Again, with the exact same steps as that for OGD with Lazy Projections. We arrive at \begin{equation}
    \boxed{\text{Regret}_T(x) \leq \frac{1}{2\eta (q - 1)}||x||_q^2 + \eta TL^2} 
\end{equation}
In particular, if all the iterations $x$ are such that $x \in S = \{x : ||x||_q \leq B\}$, and that $\eta = \frac{B}{L\sqrt{\frac{2T}{q-1}}}$, then we have that $$\text{Regret}_T(S) \leq BL\sqrt{\frac{2T}{q-1}}$$.
\end{frame}

\begin{frame}{Versatility of $p$-norm}
    Recall that $\frac{1}{p} + \frac{1}{q} = 1$. So $p \geq 2 \implies q \in (1, 2]$. Let us examine what the behaviour of the algorithm is by analysing $h(x)$ at both extremes of $q$.
    \begin{itemize}
        \item Case 1: (Easier) $q = 2 \implies h(x) = \frac{1}{2\eta}||x||_2^2$ (Euclidean Regulariser) \pause
        \item Case 2: (Tricky) $q \to 1^+$: \begin{equation*}
            h(x) = \frac{||x||_q^2}{2\eta(q-1)} \to h_1(x) = \frac{||x||_q^2 - 1}{2\eta(q-1)}
        \end{equation*}
    \end{itemize}\pause
    Why is this the case? The next iterate, $x^{(t+1)} = g(\theta^{(t+1)})$ remains the same (a.k.a algorithm doesnt change)
    \begin{align*}
    g(\theta) &= \arg\max_{x \in S} \{\langle \theta, x \rangle - h(x)\}= \arg\max_{x \in S} \left\{\langle \theta, x \rangle - \left(h(x) - \frac{1}{2\eta (q-1)}\right)\right\}\\
     &= \arg\max_{x \in S} \{\langle \theta, x \rangle - h_1(x)\}
\end{align*}
\end{frame}

\begin{frame}{Versatility of $p$-norm}
    How does $h_1$ help in our analysis?
    \begin{align*}
    \lim_{q \to 1^+} h_1(x) &= \lim_{q \to 1^+} \frac{||x||_q^2 - 1}{2\eta(q-1)}\\
    &= \lim_{q \to 1^+}\frac{\frac{\partial}{\partial q}[(\sum_{i=1}^d x_i)^{2/q}-1]}{\frac{\partial}{\partial q}(2 \eta (q-1))} \hspace{1.5cm}(\because \text{L'Hopital's Rule})\\
    &= \frac{||x||_1^2(-2 \log(\sum_{i=1}^d |x_i|)+ 2 \frac{\sum_{i=1}^d |x_i|\log|x_i|}{\sum_{i=1}^d |x_i|})}{2\eta}
\end{align*}\pause
Using the fact that $x \in \Delta_d$, we have $\sum_{i=1}^d |x_i|=1$ and $||x||_1^2 = 1$. We get 
\begin{equation*}
    \lim_{q \to 1^+} h_1(x) = \frac{1}{\eta} \sum_{i=1}^d x_i \log x_i\qquad(\text{Entropic Regulariser})
\end{equation*}
\end{frame}

\begin{frame}{Versatility of $p$-norm}
    Indeed, what we have shown is that by varying $q$ (indirectly via $p$), we are able to interpolate between the properties of the entropic and Euclidean regularisations. Furthermore, if we let $p = \log d$, we can have: \begin{align*}
    \text{Regret}_T(S) &\leq BL \sqrt{\frac{2T}{q-1}}\\
    &= BL \sqrt{\frac{2T}{1/(p-1)}} &\because \,\frac{1}{p} + \frac{1}{q} = 1\\
    &= BL \sqrt{2T(\log d - 1)} = \mathcal{O}(BL\sqrt{T\log d})
\end{align*}
\end{frame}

\begin{frame}{What we have seen}
    \begin{itemize}
        \item Brief Recap for FoReL and its regret bounds
        \item Classical Mirror Descent 
        \item Online Mirror Descent and its relation to FoReL
        \item Variants of OMD algorithms and their regret bounds
    \end{itemize}
    \pause
    We have seen how we (heavily) utilised results from the FoReL framework to derive the regret bounds of different variants of OMD algorithms. We now ask the following:
    \begin{block}{}
        Can we have a much \textit{simpler} approach where we can obtain \textit{tighter} bounds?
    \end{block}
\end{frame}
\section{Tighter Regret Bounds with Duality and Local Norms}
\begin{frame}{}
    We now shift away from using the FoReL framework to study the regret bounds of various OMD algorithms and instead adopt a new lens to perhaps get tighter bounds and in a much simpler way. In this section, we introduce two other ways to analyse regret bounds.
    \begin{itemize}
        \item Duality
        \item Local Norms
    \end{itemize}
\end{frame}

\begin{frame}{A Primer on Fenchel Conjugacy}
We first give the definition of the Fenchel conjugate of a convex function $f$:
\begin{block}{Fenchel Conjugate of a convex function $f$}
    The Fenchel conjugate of the convex function $f$, denoted as $f^*$ is defined to be\begin{equation}\label{eqn:frenchel_conjugate}
        f^*(\theta) = \sup_{x \in \text{dom}(f)} \langle \theta, x \rangle -f(x)
    \end{equation}
\end{block}
Intuitively, we can think of the Fenchel conjugate of $f$ as an alternative representation of a function.\pause
\begin{equation*}
    \text{Set of points} : (x, f(x)) \to \text{Set of tangents} : (\theta, -f^*(\theta))
\end{equation*}\pause
The Fenchel conjugate $f^*$ is the function which takes in $\theta$ and outputs the \textbf{negative} $y$-intercept of the tangent line with the slope $\theta$, $f^*(\theta)$.
    
\end{frame}

\begin{frame}{Fenchel Conjugate Pairs examples}
\begin{table}[h]
\centering
\renewcommand{\arraystretch}{1.3}
\begin{tabular}{ccl}
\toprule
$f(x)$ & $f^\star(\theta)$ & Comments \\
\midrule
$\frac{1}{2}\|x\|_2^2$ & $\frac{1}{2}\|\theta\|_2^2$ & \\
$\frac{1}{2}\|x\|_q^2$ & $\frac{1}{2}\|\theta\|_p^2$ & where $\frac{1}{p} + \frac{1}{q} = 1$ \\
$\sum_i x[i]\log(x[i]) + I_{\Delta_d}(x)$ & $\log\big(\sum_i e^{\theta[i]}\big)$ & (normalized Entropy) \\
$\frac{1}{\eta} g(x)$ & $\frac{1}{\eta} g^\star(\eta\theta)$ & where $\eta > 0$ \\
$g(x) + \langle x, x' \rangle$ & $g^\star(\theta - x')$ & \\
\bottomrule
\end{tabular}
\caption{Some useful Fenchel conjugate pairs.}
\label{tab:fenchel-pairs}
\end{table}
\end{frame}

\begin{frame}{Strong-Smoothness}
    Recall the definition of $\sigma$-strong convexity. \pause
    \begin{block}{$\sigma$-strong convexity}
        A function $f$ that is $\sigma$-strongly convex over a set $S$ with respect to norm $||\cdot ||$ if for any $x' \in S$, we have that \begin{equation*}
        \forall z \in \partial f(x'),\;\forall x \in S, \; f(x) \geq f(x') + \langle z, x - x' \rangle + \frac{\sigma}{2} || x - x' ||^2
    \end{equation*}
    \end{block}
    \pause
    A related property of strong-convexity is what we call \textit{strong-smoothness} and it is related to the Bregman divergence
    \begin{block}{$\sigma$-Strong-smoothness}
        A function $h$ is $\sigma$-strong smooth with respect to a norm $||\cdot||$ if it is differentiable and for all $x, x'$, we have that
        \begin{equation}
            D_h(x||x') \leq \frac{\sigma}{2}||x - x'||^2
        \end{equation}
    \end{block}
\end{frame}

\begin{frame}{Strong-convexity-Strong-smoothness duality}
    \begin{block}{Strong-convexity-smoothness Duality Lemma}
        Assume that $h$ is a closed and convex function. $h$ is $\sigma$-strongly convex with respect to the norm $||\cdot||$ if and only if the Fenchel conjugate $h^*$ is $\frac{1}{\sigma}$-strongly smooth with respect to the dual norm $||\cdot||_*$.
    \end{block}
\end{frame}
\end{document}